% ====== Paquetes básicos ======
\usepackage[utf8]{inputenc}    % Codificación UTF-8
\usepackage[T1]{fontenc}       % Acentos correctos
\usepackage[spanish]{babel}    % Traducción al español
\usepackage{csquotes}          % Recomendado para biblatex con babel
\usepackage{hyperref}          % Hipervínculos
\usepackage{amsmath,amssymb}   % Matemáticas
\usepackage{graphicx}          % Imágenes
\usepackage{tikz}              % Diagramas vectoriales
\usetikzlibrary{arrows.meta,positioning}
\usepackage{caption}           % Mejor control de captions
\usepackage{geometry}          % Márgenes
\usepackage{fancyhdr}          % Encabezados y pies
\usepackage{parskip}           % Mejor separación entre párrafos

% ====== Referencias ======
\usepackage[
    backend=biber,
    style=numeric-comp,     % cambie a APA de ser necesario
    sorting=none
]{biblatex} 
\bibliography{references}

% agrege `\printbibliography{}` para las referencias
% cite con `\cite{}`
% agregue las referencias a un `references.bib`

% ====== Configuración ======
\geometry{margin=2.5cm}
\pagestyle{fancy}
\fancyhf{}
\rhead{\thepage}
\setlength{\headheight}{13pt}
\lhead{\leftmark}

% ====== Comando personalizado para imágenes ======
% Uso: \imagen{ruta}{Texto del pie de figura}
\newcommand{\imagen}[2]{
    \begin{figure}[h!]
        \centering
        \includegraphics[width=200px]{#1}
        \caption{#2}
    \end{figure}
}

% ====== Datos del artículo ======
\title{Portada}
\author{Lester Iván Hernández Hernandez}
